\chapter*{Introduction}

Dans ce rapport nous étudions l'équation Benjamin-Bona-Mahony, un modèle 
pour les ondes de surface de grande longueur devant leur amplitude. Ce modèle est une variante 
d'une équation bien connue parmi les équations aux dérivées partielles : l'équation Korteweg - de Vries (KdV). Cette équation,
parue à la fin du $\text{XIX}^e$ siècle, avait pour objectif initial de modéliser des vagues dans un milieu 
particulier (non-linéaire et dispersif) en faible profondeur (typiquement dans un canal). Cette équation a la particularité
d'admettre une solution étonnante : le soliton. Il s'agît d'une vague capable de se propager sur de très longues distances sans toutefois que son
profil ne varie. Cette dernière a la particularité de pouvoir interagir avec d'autres soliton tout en 
reprenant sa forme initiale (là où d'autres vagues interagiraient de manière chaotique à première vue).
Cette équation a suscité beaucoup d'intérêt à partir des années 60, où des avancées ont été faites, numériquement, sur
le comportement asymptotique d'une perturbation selon ce modèle. Ces avancées ont ensuite été vérifiées théoriquement. 
Ce modèle a reçu beaucoup d'attention dans les années qui ont suivi (notamment pour l'étude du caractère "bien posé" du problème), puis 
d'autres domaines de la physique ont pu s'approprier ce modèle. Nous pouvons par exemple retrouver l'équation KdV dans le contexte de 
vagues hydromagnétiques de plasma sans collision, de longues vagues dans les cristaux non harmoniques, pour les plasmas ion-acoustique ou encore en cosmologie.

En 1972, trois mathématiciens anglais et américains proposent une alternative au modèle, l'équation qui portera leurs nom conjoints, Benjamin-Bona-Mahony. Cette proposition
repose sur l'idée d'améliorer la stabilité numérique du modèle initial KdV, en particulier pour des ondes plus courtes (là où le modèle KdV atteingnait un limite, en raison de sa relation de dispersion non bornée).
Dans leur article initial, les trois mathématiciens dérivent leur modèle tout en le comparant à KdV et démontrent les principales propriétés
mathématiques associées.

Dans ce rapport nous reprendrons les principaux résultats de leur article tout en poussant l'étude sous l'angle de l'analyse numérique, à l'aide de divers schémas numériques, dont nous expliciterons quelques propriétés essentielles.
Notre travail repose sur la compilation de savoirs déjà existants sur les différents sujets que nous abordons. Notre rapport s'inscrit donc, en toute humilité, dans une démarche d'entretient des savoirs
existants.
Bien que ce rapport ne constitue qu'une approche primaire du sujet, nous avons su apprécier cette ouverture vers la recherche sur le thème de l'océanographie, l'analyse numérique ou encore l'analyse fonctionnelle.

\vspace{1em}

Nous souhaitons remercier notre encadrant, Fabien Marche, qui a su nous stimuler à la hauteur de l'intérêt que nous portions sur le sujet, en prenant, notamment, le temps et le soin de nous proposer différentes pistes d'étude 
tout au long de ce projet. Nous souhaitons également le remercier pour son humanité et les riches discussions que nous avons pu avoir avec lui, non pas simplement sur les mathématiques, mais sur le contexte sociétal et sur l'enjeu de ces dernières de nos jours.
Nous espérons que ce rapport sera à la hauteur de l'engagement dont nous avons fait preuve durant ces quelques mois de recherches sur ce sujet.